\chapter*{Preface}
\paragraph{What is Commutative Algebra?}
The origin of commutative algebra comes from the following
two fields:
\begin{itemize}
\item \textbf{algebraic number theory,}
which studies prime factorization in number fields (such as $\mathbb Q(\sqrt{-5})$)
and its ring of integers (such as $\mathbb Z$ and $\mathbb Z[\sqrt{-5}]$).
\item \textbf{algebraic geometry,} which studies 
solutions of systems of algebraic equations 
(such as $x^2+y^2=1$), which can be turned 
to study of quotients of polynomial rings 
(such as $\CC[x,y]/(x^2+y^2-1)$).
\end{itemize}
Throughout time, people have realize that theory from these 
two subjects can be unified into a single theory of 
commutative rings. This results in a single, clean theory of commutative algebra that stands on its own, but its true purpose is to be applied in the context of number theory and algebraic geometry
\paragraph{Prerequisite}
We assume that the reader is familiar with rings and modules at an undergraduate level (e.g., as in \cite{artin}). 
Basic concepts of rings and modules will be rapidly reviewed 
in \Cref{chap:rings} and \Cref{chap:modules}.
Familiarity with basic category theory and topology 
is assumed, but we quickly review those in \Cref{chap:background}.
Knowledge of field theory at an undergraduate level is assumed in some places. 

\paragraph{Content}
These notes cover standard material typically covered in an introductory graduate course (such as 18.705 at MIT).
The goal is to develop basic theory of rings and modules
necessary in the study of algebraic number theory 
and algebraic geometry. 
Although we do not cover algebraic geometry
or algebraic number theory,
we try, whenever possible, to point out motivations and 
connections from these two subjects.
The breakdown by chapter is as follows.
\begin{itemize}
    \item \textbf{Chapter 1 (Rings).} 
    We start out by quickly reviewing the notions of rings and ideals, including primary and maximal ideals.
    The spectrum of a ring, which is a central notion in algebraic geometry, is introduced in \Cref{sec:spec}.
    \item \textbf{Chapter 2 (Modules).}
    We quickly review basis notions of modules,
    then we proceed to discuss properties and constructions related to modules at a significantly greater depth
    than typically covered in undergraduate algebra.
    We discuss Cayley-Hamilton theorem from linear algebra and its consequence on modules that we will use multiple times throughout.
    Then, we discuss important constructions related to modules: exactness, colimits, and tensor products.
    \item \textbf{Chapter 3 (Localization).}
    Localization is the ``art of adding denominators'' to rings. 
    This construction has a geometric interpretation of ``zooming into'' open sets or points
    and is immensely useful in studying rings.
    \item \textbf{Chapter 4 (Integral Extensions).}
    Integral extension and integral closure is a fundamental notion that connects to finite maps, smooth curves, and ring of algebraic integers.
    As a consequence of this theory, we present Noether normalization lemma and use it to prove Hilbert nullstellensatz and study dimension of a variety.
    \item \textbf{Chapter 5 (Chain Conditions).}
    We discuss Noetherian rings and modules, the fundamental finiteness conditions that has numerous consequences.
    Artinian rings are also discussed.
    \item \textbf{Chapter 6 (Associated Primes and 
    Primary Decomposition).}
    In unique factorization domains such as $\ZZ$,
    every element can be factored uniquely into products of primes.
    We discuss a analogue of this in a general ring,
    which will allow us to decompose an ideal into intersection of primary ideals (which is an analogue of prime power).
    \item \textbf{Chapter 7 (DVR and Dedekind Domains).}
    Discrete valuation rings (DVR) are the simplest kind of rings after fields. 
    Geometrically, its ``valuation'' comes naturally from trying to define order of vanishing of functions at a point.
    A global analogue of DVR is a Dedekind domain,
    which arise naturally from ring of integers (in number theory) and a smooth curve (in algebraic geometry).
    We discuss their properties.
    \item \textbf{Chapter 8 (Completion).}
    As power series is useful in studying functions in real and complex analysis, we develop the notion of power series 
    in commutative algebra, resulting in construction known as completion.
    We prove several nice properties of completion.
    \item \textbf{Chapter 9 (Dimension Theory).}
    Although dimension is an intuitive concept in geometry,
    proving properties about dimension of rings are notoriously difficult.
    In this chapter, we develop a robust framework to study and compute dimensions.
    A brief treatment of regular local rings
    (rings corresponding to smooth varieties) is discussed.
    \item \textbf{Chapter 10 (Flatness).}
    We study flat modules (defined in Chapter 2 in terms of tensor products) in more detail, discussing various geometric consequence such as going-down property, and constant dimension of fiber.
\end{itemize}
The text is meant to be read in order, although dependencies between many chapters are weak, and some flexibility is possible.

All contents existed in standard commutative algebra textbooks.
The one closest to our notes is \cite{atiyah_macdonald}.
Other books that influenced our presentation are \cite{altman_kleiman}, \cite{eisenbud}, and \cite{vakil}.
I claim no originality in any of these materials.

\paragraph{History and Acknowledgements}
These notes were based on the notes that I took when taking 
the Fall 2023 instance of 18.705 (a commutative algebra course) at MIT,
taught by Davesh Maulik.
Since then, many parts were rewritten, and a new \Cref{chap:flatness} on flatness was added.
Nevertheless, the structure of these notes still largely follows how Maulik taught the class, and his insights and presentation is still imparted throughout these notes.
The flatness chapter was based on the Fall 2024 instance of 18.705, taught by Miguel Moreira, which also influences the presentation of various parts of theses notes as well.

We thank Derek Liu for contributing part of \Cref{chap:flatness}.
We also thank Srinivas Arun, David Dong, and Linus Tang
 for useful comments and suggestions.
